\chapter{数字市场竞争与平台反垄断}

\begin{examguide}
\textbf{本章考情}：平台反垄断是 0258 论述题与案例题的最高频命题领域（\important{专硕·核心}），SSNIP 测试与市场势力测度是计算题的命题素材；学硕微观中作为垄断理论的前沿应用（\important{学硕·热点}）。

\textbf{核心考点}：（1）数字市场垄断的形成机制（数据反馈回路、守门人）；（2）新型垄断行为与标志性案例（自我优待、杀手型并购、数据驱动型并购、二选一、算法合谋）；（3）市场界定专题：SSNIP 测试、临界损失分析、玻璃纸谬误与多边市场界定；（4）HHI 与勒纳指数；（5）中国平台反垄断制度与标志性案例（含支配地位推定条款）；（6）欧盟 DMA 的事前监管逻辑；（7）App Store 抽成之争与 Epic 案的市场界定范式。

\textbf{本章主线}：垄断如何形成（6.1）→ 新型垄断行为（6.2）→ 如何界定市场与测度势力（6.3）→ 制度与实践（6.4）。
\end{examguide}

\section{数字市场垄断的形成}

\subsection{三重机制的叠加}

数字市场垄断的形成机制可以概括为\important{三重力量的叠加}：\textbf{网络效应}（第3章）使大平台的产品价值更高，\textbf{规模经济}（第1章零边际成本）使大平台的成本更低，\textbf{数据壁垒}（data barrier）使大平台的产品质量更高且壁垒随时间加厚。前两者在传统产业中同样存在，数据壁垒则是数字市场的特有机制，其作用机制可以用一个简单的动态模型刻画。

\begin{derivation}{数据反馈回路}{}
\textbf{数据反馈回路的增长条件。}设平台在 $t$ 期的用户数据存量为 $d_t$，数据提升服务质量（推荐精度、风控能力），服务质量吸引用户，用户产生更多数据，这一数据反馈回路（data feedback loop）可形式化为
\[
d_{t+1} = f(d_t), \qquad f(0) = 0,
\]
其中 $f$ 递增（$f' > 0$：数据越多，下期产生的数据越多），且数据转化的边际贡献递减（$f'' < 0$：推荐精度的边际提升随存量下降）。$f'(d)$ 的直观含义是数据的边际"繁殖率"：一单位新增数据在下一期再产生多少新数据。两个条件刻画数据壁垒的成与败。其一，\textbf{排他性}：用户数据不跨平台共享，在位者的 $f$ 不能为进入者所用；其二，\textbf{启动条件} $f'(0) > 1$：小数据量时边际繁殖率大于 1，微小的数据优势会被放大而非衰减。

\textbf{进入者的死结。}由 $f(0) = 0$，$d = 0$ 是一个\textbf{吸收不动点}：无数据者永远无数据。进入者的初始存量恰为 $d_0 = 0$，其服务质量起点系统性落后于在位者，而用户迁移又需要服务质量的跨越式提升：\textbf{无数据则无质量、无质量则无用户、无用户则无数据}，数据壁垒使进入者被锁死在循环之外。

\textbf{壁垒的饱和点。}当边际繁殖率随存量衰减至 1 以下时（$\lim_{d \to \infty} f'(d) < 1$），方程 $d^{*} = f(d^{*})$ 存在唯一的正解：$f$ 与 45° 线的唯一正交点 $d^{*}$ 是稳定稳态（在该点 $f'(d^{*}) < 1$，扰动会被循环吸收回 $d^{*}$）。\important{数据壁垒的长期高度就是 $d^{*}$}，它取决于数据转化的边际效率（$f$ 的形状）而非时间累积：转化效率高的在位者壁垒更高，但任何在位者的数据优势都存在天然的饱和点。
\end{derivation}

\begin{figure}[htbp]
\centering
\begin{tikzpicture}[>=Stealth,
    node/.style={rectangle, rounded corners=2.5mm, draw=second, thick, fill=second!8, align=center, font=\small, minimum width=26mm, minimum height=12mm},
    arr/.style={->, very thick, main}]
% 节点先定义坐标（主层），连线放入背景层，节点边框盖住线头
% 标签规则：斜边标签沿箭头方向紧贴平行（sloped, above）；竖边标签垂直贯穿箭头（白底骑线）
\node[node] (U) at (0,2.4) {用户规模\\$n \uparrow$};
\node[node] (D) at (4.2,0) {数据存量\\$d \uparrow$};
\node[node] (Q) at (0,-2.4) {服务质量\\$q(d) \uparrow$};
\begin{scope}[on background layer]
\draw[arr, shorten <=-2pt, shorten >=-2pt] (U.east) -- node[font=\small, color=structurecolor, sloped, above, pos=0.5] {使用与互动} (D.north west);
\draw[arr, shorten <=-2pt, shorten >=-2pt] (D.south west) -- node[font=\small, color=structurecolor, sloped, below, allow upside down=false, pos=0.5] {训练与优化} (Q.east);
\draw[arr, shorten <=-2pt, shorten >=-2pt] (Q.north) -- node[font=\small, color=structurecolor, fill=white, inner sep=1.5pt, pos=0.5] {吸引与留存} (U.south);
\end{scope}
\end{tikzpicture}
\caption{数据反馈回路}
\label{fig:dataloop}
\end{figure}

图~\ref{fig:dataloop} 的循环刻画了数据壁垒的形成机理，三条边分别是反馈回路的三个传递环节：用户使用产生数据（使用与互动），数据训练改进服务（训练与优化），服务质量吸引更多用户（吸引与留存）。形式化结果给出两个政策落点。其一，进入者被困在吸收不动点 $d = 0$，反垄断政策的回应选项都旨在为进入者提供循环的"替代起点"：数据互操作（interoperability）与数据可携带（data portability）提高进入者的初始存量，并购审查则约束在位者对潜在威胁的收购。其二，壁垒的长期高度是稳态 $d^{*}$ 而非无限累积，规制重点是数据转化的边际效率环节而非存量本身。

\subsection{超级平台与守门人}

当平台在多个相邻市场拥有市场支配地位（market dominance）时，其性质发生质变：从单一市场的经营者变为多个市场的\textbf{基础设施}。最直观的例子是应用商店：开发者要接触 iPhone 用户，除苹果 App Store 外别无通道；搜索引擎之于网站、操作系统之于应用，同样构成这样的"必经之路"。欧盟《数字市场法》（Digital Markets Act，DMA）用\textbf{守门人}（gatekeeper）称呼这类平台：它们像体育场的大门，观众进场必须经过它。

\begin{definition}[守门人（gatekeeper）]
守门人是指在核心平台服务（搜索引擎、社交网络、应用商店、操作系统等）中具有稳固且持久的地位、并对商业用户接触最终用户构成控制性瓶颈的平台。
\end{definition}

定义中的三个限定词各对应一个判定要素，逐一拆开理解。"\textbf{核心平台服务}"指搜索引擎、社交网络、应用商店、操作系统等数字经济的枢纽业务，守住枢纽，商业用户就绕不开平台。"\textbf{稳固且持久}"要求优势不是暂时的：DMA 以连续三年达到用户规模门槛作为客观标准，避免执法机关逐年论证支配地位。"\textbf{控制性瓶颈}"指商业用户接触最终用户必须经过平台：开发者离开 App Store 就接触不到 iOS 用户，网站离开搜索引擎就流失大半流量。

守门人概念的价值不在称谓，而在它改变了监管方式。传统反垄断的做法是"出事了再查"：先证明平台有支配地位、实施了滥用行为、造成了竞争损害，再处罚，谷歌购物案的处理就是如此。守门人制度则反过来：先把"大到绕不开"的平台直接列出来，再事先规定它们能做什么、不能做什么。为什么需要这种转变？因为数字市场的正反馈（第 6.1.1 节）使竞争损害一旦形成就难以逆转，等事后查处时市场往往已收敛于单一平台，处罚换不回竞争。\mn{守门人三门槛口诀：影响大（营业额/市值）、用户多（活跃用户）、地位稳（连续三年）。}守门人的边界有两点：其一，守门人是\important{结构认定}，取决于规模与地位的客观状态，与是否实施了违法行为无关；其二，对守门人的规制是\important{行为义务}而非结构拆分，义务清单针对特定行为设定（第 6.4 节）。

\section{新型垄断行为}

垄断地位一旦形成，平台便具备了实施新型垄断行为的条件。本节考察四类数字市场特有的行为：自我优待、杀手型并购、"二选一"与算法合谋。前三类分别是平台对裁判权、未来竞争与用户锁定的利用，算法合谋则是数字技术带来的全新挑战。四类行为各有标志性执法先例，也各有独立的经济分析，本节逐一展开。

\subsection{自我优待}

传统滥用行为分析的前提是两个相互独立的经营者之间的市场交易；数字平台的结构创造了新的行为类型：平台既控制上游的基础设施（搜索引擎、应用商店、操作系统），又在下游市场上与依赖该基础设施的独立经营者竞争。这种双重身份使平台有能力也有激励偏向自己的产品，即自我优待。

\begin{definition}[自我优待（self-preferencing）]
自我优待指平台在自己的基础设施上为自己的产品提供优于竞争对手的待遇：搜索结果中优先展示自营商品、应用商店中优先推荐自有应用、操作系统预装自有服务。
\end{definition}

自我优待的竞争损害经由纵向结构传导：平台在基础设施环节的支配地位，通过排序、推荐、预装等流量分配机制传导至下游市场。设下游市场上存在平台自营产品 $A_P$ 与独立产品 $A_I$，两者初始质量相当（$q_P = q_I$），竞争损害可以分解为两个效应，形式化如下。

\begin{derivation}{自我优待的竞争损害}{}
\textbf{需求转移效应。}设平台总流量为 $N$，每单位流量的价值为 $v$。自我优待使自营产品的发现概率提高 $\Delta r$，独立产品的流量相应下降 $\Delta r \cdot N$，由此转移的利润为 $\Delta r \cdot N \cdot v$。这一效应是静态的：它不改变产品质量，只改变流量的分配。

\textbf{投资抑制效应。}独立产品以质量投资 $e$ 应对竞争，投资通过用户反馈转化为发现概率：$r_I = (\bar r - \Delta r)\, h(e)$，其中 $\bar r$ 为无优待时的发现概率，$h' > 0$（投资提升质量与口碑）、$h'' < 0$（边际回报递减），投资成本为 $c(e)$（$c' > 0$、$c'' > 0$）。乘性设定刻画的是"投资的效果需要通过用户使用来实现"：流量被截留时，同样的投资改善触及的用户更少，这与 6.1 节的数据反馈回路同构。独立产品的利润为
\[
\pi_I(e) = (\bar r - \Delta r)\, h(e)\, N v - c(e),
\]
一阶条件（投资的边际收益等于边际成本）：
\[
(\bar r - \Delta r)\, h'(e^*)\, N v = c'(e^*).
\]
对 $\Delta r$ 隐函数求导：
\[
\frac{de^*}{d\Delta r} = \frac{h'(e^*) N v}{(\bar r - \Delta r) h''(e^*) N v - c''(e^*)} < 0,
\]
分母两项均为负（$h'' < 0$、$-c'' < 0$），分子为正，故\important{自我优待程度上升，独立产品的均衡质量投资下降}。
\end{derivation}

两个效应叠加的长期后果是"自我实现的劣化"：投资抑制使 $q_I$ 相对 $q_P$ 下降，质量差距又成为平台继续优待的"证据"——优待制造了优待的理由，这一循环与 6.1 节的数据反馈回路同构。\mn{辨析：自我优待本身不当然违法，违法的是"实质性排除竞争且无效率抗辩"的优待（谷歌购物案逻辑）。}竞争损害的实质在于：平台\important{既是运动员又是裁判员}，其在基础设施环节的垄断势力被传导至下游市场。规制边界在于区分"合法的差异化经营"与"滥用基础设施势力"：当优待行为实质性削弱下游竞争且缺乏效率抗辩时，构成滥用市场支配地位（abuse of dominance）（谷歌购物案的处理逻辑）。

自我优待的三类行为各有执法先例。\textbf{搜索结果优待}：欧盟 2017 年认定谷歌将自营比价购物服务置顶、系统降权独立比价网站，独立比价网站的流量断崖式下跌（部分网站点击量下跌超过八成），处罚 24.2 亿欧元（第 6.4 节案例）。\textbf{应用商店优待}：意大利竞争管理局 2021 年认定亚马逊在自营物流（FBA）与"购买按钮"（Buy Box）上给予自营商品特权、排斥独立卖家，罚款 11.3 亿欧元；美国联邦贸易委员会 2023 年就同类行为起诉亚马逊。\textbf{操作系统预装}：欧盟 2018 年认定谷歌利用安卓系统的预装与默认设置绑定自有搜索与浏览器，处罚 43.4 亿欧元。三个案子的共性在于：优待的对象都是\textbf{自家产品}，手段都是\textbf{基础设施层面的规则设定}，损害的论证都落在\textbf{下游竞争被实质性削弱}——这也是答题时把本节推导（需求转移 + 投资抑制）翻译成案例语言的三个抓手。

\subsection{杀手型并购}

传统并购审查关注的是横向合并：两家在位企业合并，市场份额叠加。数字市场出现了一类不同的并购：在位平台收购市场份额尚小、但掌握前沿技术与早期用户的初创企业，其动机可能不是整合其产品，而是消除其未来的竞争威胁，这类并购被称为杀手型并购。

\begin{definition}[杀手型并购（killer acquisitions）]
杀手型并购指在位平台收购尚处早期、但未来可能构成竞争威胁的创新企业，目的不是整合其产品而是\textbf{扼杀}其产品线。
\end{definition}

杀手型并购的规制难点在于其威胁的形态是\textbf{未来的竞争（future competition）}而非当下的市场份额：被收购的初创企业尚未形成规模，传统并购申报门槛（营业额）经常无法触发审查，形成监管失灵的制度性漏洞。\mn{杀手型并购的识别难题：初创企业营业额小，常规申报门槛常无法触发审查；欧盟《数字市场法》第 14 条要求守门人申报拟议并购（无营业额门槛，"抢跑"条款），美国 2021 年撤销《纵向并购指南》收紧数字并购审查。}各国回应的方向一致：以交易对价而非营业额为审查门槛，对数字并购从严审查。欧盟《数字市场法》第 14 条要求守门人就拟议并购向欧委会申报，不设营业额门槛；中国《反垄断法》第 26 条规定，未达申报标准但有证据证明集中具有或可能具有排除、限制竞争效果的，执法机构可以要求申报。并购控制的福利标准在数字市场面临 Williamson 权衡（Williamson tradeoff）的变形。

\begin{derivation}{Williamson 权衡}{}
\textbf{Williamson 权衡。}并购的竞争效应由两个方向相反的力构成：并购可能产生成本效率（研发整合、规模经济），使边际成本从 $c_1$ 降至 $c_2$；同时并购增强市场势力，使价格从 $p_1$ 升至 $p_2$。设并购前市场为竞争性定价 $p_1 = c_1$，$p(q)$ 为反需求函数，$q(p_2)$ 为并购后产量（$q(p_2) < q(p_1)$）。社会福利的净变化为
\[
\Delta W = \underbrace{\int_{c_2}^{c_1} q(p_2)\, dc}_{\text{成本效率收益}} \;-\; \underbrace{\int_{q(p_2)}^{q(p_1)} \bigl[p(q) - p_1\bigr]\, dq}_{\text{价格上升的无谓损失}} .
\]
第一项是\textbf{成本节约矩形}：成本下降 $c_1 - c_2$ 只发生在并购后仍产出的 $q(p_2)$ 个单位上，故沿成本轴积分时被积函数固定为 $q(p_2)$，即 $(c_1 - c_2)\, q(p_2)$。第二项是\textbf{无谓损失三角}：产量从 $q(p_1)$ 收缩至 $q(p_2)$，损失的社会剩余以反需求曲线为界、以并购前成本线 $p_1$ 为底：被放弃的每个单位对消费者的价值 $p(q)$ 超出其并购前生产成本 $p_1$。并购增进总福利当且仅当矩形超过三角。
\end{derivation}

传统并购审查中，成本效率收益与无谓损失大体可比，净效应取决于具体参数。数字市场中这一权衡被系统性打破：零边际成本产业中成本效率收益的空间极小（$c_1$ 已近于零，$c_1 - c_2$ 有限），而网络效应使市场势力增强的持续性远超传统产业，并购消除的不仅是在位竞争，更是\important{未来的潜在竞争}。因此数字并购审查需要更严格的标准：能够触发审查的门槛应更低（关注交易对价而非营业额），抗辩的效率主张需要更高的证明标准。

杀手型并购最著名的例证是 Facebook 对 Instagram（2012 年，约 10 亿美元）与 WhatsApp（2014 年，约 190 亿美元）的收购：被收购时 Instagram 仅有 13 名员工、尚无收入，WhatsApp 几乎不产生营业额，传统申报门槛在二者身上完全失效。美国联邦贸易委员会 2020 年 12 月起诉 Facebook，主张这两笔收购正是为消除潜在竞争威胁而实施、要求拆分；2021 年 6 月法院以市场支配地位论证不足为由驳回起诉，联邦贸易委员会于 2021 年 8 月修改诉讼状重新起诉，诉讼至今悬而未决。案件的关键在于证明反事实竞争：如果没有收购，Instagram 与 WhatsApp 本会成为 Facebook 的竞争者。这一证明困难是杀手型并购审查的普遍难题：执法者需要论证的是一项尚未发生的竞争。实证研究（Cunningham–Ederer–Ma，2021）在制药业找到了系统证据：被在位企业收购的候选新药，其研发项目被终止的概率显著更高，且终止的恰恰是与在位者产品线重叠的项目，"扼杀"假说由此获得数据支持。

与杀手型并购并列的是\textbf{数据驱动型并购}（data-driven mergers）：并购的目标不是产品线或初创团队，而是数据资产本身。收购方以交易获得目标企业的用户数据、行为记录与算法积累，将其并入自身的数据反馈回路（6.1 节），在并购审查未触及的维度上扩大数据壁垒。数据驱动型并购的竞争损害有三：其一，\textbf{数据封锁}，竞争者失去该数据源的接入可能，数据市场的进入成本被抬高；其二，\textbf{创新转移}，目标企业的数据资产与产品路线被吸收后，独立的创新方向消失；其三，\textbf{审查逃逸}，数据资产不体现为营业额，传统申报门槛对"为数据而并购"的交易系统性失灵。\mn{数据驱动型并购与杀手型并购的区别：前者目标是数据资产与反馈回路，后者目标是消除未来竞争的产品线；两者常同时发生，审查难点都在"营业额门槛失灵"。}中国《反垄断法》2022 年修订后第 26 条第 2 款授权对未达申报标准但有排除限制竞争效果的集中进行调查，其立法动机之一正是覆盖此类并购，执法实践与理论工具的同步推进构成数字并购治理的主线。

\subsection{"二选一"与算法合谋}

\textbf{"二选一"}（排他性交易）的竞争分析在第4章竞争瓶颈模型的基础上推进：平台利用对商家的锁定，强制商家放弃竞争对手平台，其效果是\important{抬高竞争对手的商家获取成本，保护甚至扩张自身的单归属商家池}，进而为商家侧的垄断收费提供制度保障。第4章给出的定性判断变量——平台市场势力的大小、多归属选择是否被实质剥夺——在这里获得执法应用：中国市场监管总局对阿里巴巴"二选一"行为的处罚确立了处理框架（第 6.4 节阿里巴巴案）。执法披露的细节量化了锁定的规模：美团案中，2018 至 2020 年签订独家协议的商家约 163 万家，累计收取独家合作保证金约 12.89 亿元，非独家商家的佣金率高出约 5 至 7 个百分点（国家市场监督管理总局处罚决定书，2021），价格差异正是竞争瓶颈结构在商家侧收费上的直接体现。阿里巴巴案将相关市场界定为"网络零售平台服务市场"，美团案则界定为"网络餐饮外卖平台服务市场"：同一行为在不同案件中因市场界定不同而有不同的支配地位论证，市场界定先行的实践逻辑在第 6.3 节展开。

\textbf{算法合谋}（algorithmic collusion）是数字市场的全新挑战。传统合谋需要人为的意思联络，算法使默示合谋（tacit collusion）的达成条件发生质变。第2章触发策略的结论是合作需要 $\delta \geq \delta^*$，即惩罚的即时性决定合谋的可持续性；这一结论可以直接翻译为反应时间的形式化条件。

\begin{derivation}{惩罚时滞与合谋稳定性}{}
\textbf{惩罚时滞与合谋稳定性。}设利率为 $r$（连续复利），厂商观察到对方降价并实施惩罚性降价的时滞为 $\Delta t$，时滞期间的贴现因子为
\[
\delta = e^{-r\,\Delta t}.
\]
第2章的合谋条件 $\delta \geq \delta^*$ 转化为时滞条件：
\[
\Delta t \leq -\frac{\ln \delta^*}{r}.
\]
右侧为正（$\delta^* < 1$），含义是：惩罚必须在临界时滞 $-\ln \delta^*/r$ 内到来，合谋才能维持。时滞的破坏作用体现在两个方面。其一，\textbf{背叛窗口}：时滞内背叛者按接近垄断的价格获利，背叛收益 $\int_0^{\Delta t} \pi^D e^{-rs} ds \approx \pi^D \Delta t$ 与 $\Delta t$ 成正比；其二，\textbf{惩罚贴现}：惩罚损失要到 $\Delta t$ 之后才开始，折现后更小。比较两种情形：\textbf{人工定价}下，发现对手降价、召集会议、调整价格目录的周期以周计（$\Delta t$ 为数周），背叛窗口长达数周，当 $\delta^*$ 接近 1（背叛诱惑大、惩罚损失小）时条件 $\Delta t \leq -\ln \delta^*/r$ 可能被违反；\textbf{算法定价}下，定价算法以秒级频率抓取对手价格并自动响应，$\Delta t \to 0$，条件对任意 $\delta^* < 1$ 恒成立——\important{超竞争价格无需任何沟通、协议或人的共同计划，仅凭算法的即时报复机制即可自我维持}。
\end{derivation}

两种情形的对比说明，算法从两个方向拆除了人工合谋的障碍。其一，\important{报复即时化}：$\Delta t \to 0$ 使背叛窗口收缩至秒级，合谋可持续性条件对任意 $\delta^* < 1$ 恒成立。其二，\important{合谋价格的确立不再需要沟通}：人工定价下超竞争价格的确立依赖卡特尔会议之类的明示协议（这正是传统市场需要明示协议的根源），算法定价则以高频试错自动收敛到合谋价格。\mn{易错点：算法合谋的执法难点不在"合谋"的证明，而在"协议"的证明：无意思联络的均衡结果难以落入反垄断法"垄断协议"的构成要件。}算法合谋的执法难题由此产生："合谋"是算法策略收敛的均衡结果而非人的共同计划，反垄断法以"协议"为责任基础的传统框架出现结构性失灵，这是各国反垄断机构（OECD 2017 专题报告）面对的前沿难题。

执法实践印证了这一难题的分层：已经查处的算法案件都属于\important{算法执行人的协议}：美国司法部 2015 年诉 David Topkins 案（亚马逊海报销售商与他人约定使用同一算法协调价格）与英国竞争与市场管理局 2016 年对两家在线海报零售商（Trod、GB eye）的处罚，共同点是算法只是执行工具，"协议"仍然存在。真正由算法自主收敛形成的合谋尚无执法先例。\mn{算法合谋四类：信使—轴辐—可预测代理—自主机器，人的参与递减、执法难度递增。}学理上，Ezrachi 与 Stucke（2016）按人的参与程度把算法合谋分为四类：\textbf{信使}（算法执行人的协议）、\textbf{轴辐}（多个商家使用同一算法供应商的定价服务，形成事实上的价格协调）、\textbf{可预测代理}（算法各自独立学习，但都学会了"以牙还牙"的报复策略，使默示合谋稳定化）、\textbf{自主机器}（算法自发收敛到超竞争价格）。四类的责任归属难度依次递增：信使类有明确协议，可依现行法查处；自主机器类无协议、无意图，现行法框架完全失灵。

\section{相关市场界定与市场势力测度}

新型垄断行为要受到规制，执法者首先必须回答两个技术性问题：涉案行为发生在哪个"市场"、行为者的"势力"有多大。前者划定竞争损害分析的边界，后者衡量将价格定在竞争水平之上的能力，即市场势力（market power）。两者构成滥用行为认定的前提。

\subsection{SSNIP 测试与临界损失分析}

反垄断分析的第一步是界定相关市场（relevant market）：市场界定过窄会夸大行为者的势力，过宽会稀释竞争约束，两种偏差都会直接改变案件结论。标准工具是假想垄断者测试。

\begin{definition}[假想垄断者测试（SSNIP 测试）]
假想垄断者测试是界定相关市场的标准方法：假想目标产品由唯一垄断者经营，考察其能否实施"小幅但显著且非暂时性的涨价"（Small but Significant and Non-transitory Increase in Price，通常为 5\%–10\%，持续一年）而仍有利可图。
\end{definition}

若能，目标产品构成独立的相关市场；若不能（消费者转向替代品使涨价无利可图），则将最接近的替代品纳入市场再行测试，直至涨价有利可图为止。SSNIP 的思想内核是\important{以替代性划定市场边界}：消费者愿意转向哪些产品，这些产品就与目标产品处于同一市场。市场不是产品的物理归类，而是竞争约束的范围。替代性有两个考察方向：\textbf{需求替代}（消费者转向其他产品）是主要标准，\textbf{供给替代}（其他生产者能迅速转产该产品）作为补充，各国执法实践以需求替代为主。涨价是否有利可图，可以由盈亏平衡条件精确计算，这就是临界损失分析（critical loss analysis）。

\begin{derivation}{临界损失分析}{}
\textbf{临界损失分析。}设产品当前价格为 $p$，边际成本为 $c$，加价为 $m = p - c$，提价幅度为 $x$（如 $x = 5\%$）。提价的直接效果是每单位利润增加 $xp$，代价是销量损失 $\Delta q$。临界损失 $L_c$ 定义为使提价恰好不亏不赚的销量损失比例：
\[
xp \cdot (q - \Delta q) - m \cdot \Delta q = 0.
\]
令 $L = \Delta q / q$，解得
\[
L_c = \frac{x}{x + m/p}.
\]
若实际损失率 $L_{\text{实际}} > L_c$，涨价无利可图，市场界定需扩大；若 $L_{\text{实际}} < L_c$，涨价有利可图，目标产品构成独立市场。注意加价率 $m/p$ 出现在分母：\important{加价率越高的产品，临界损失越小，越容易构成独立市场}：高加价产品每损失一单位销量所放弃的利润更多，小幅销量损失就足以抵消涨价收益。
\end{derivation}

一个算例可以直观把握临界损失的规模：若 $p = 100$、$c = 60$（加价率 $m/p = 0.4$），涨价 5\% 时 $L_c = 0.05/(0.05 + 0.4) \approx 11.1\%$：实际销量损失超过 11.1\%，5\% 涨价便无利可图，市场界定需要扩大。

SSNIP 的运用有一个著名陷阱：\textbf{玻璃纸谬误}（cellophane fallacy）。1956 年美国诉杜邦公司案中，法院以杜邦玻璃纸与蜡纸、铝箔之间存在激烈替代为由，认定杜邦不具垄断地位。批评者指出，法院测量的替代性发生在杜邦\textbf{垄断价格}上：垄断者已经把价格定到需求弹性大于 1 的水平，此时任何进一步涨价都会招致大量流失，测试必然得出"市场很宽、没有垄断"的结论。玻璃纸谬误的教训是：\important{SSNIP 应从竞争性价格而非现行价格出发}，否则会在真垄断者身上测出宽市场，把垄断者放走。

SSNIP 测试在数字市场面临根本性困难，首先是\textbf{零价格问题}：搜索引擎、社交网络对消费者免费（$p = 0$），以涨价为基础的测试无从谈起。学理上的替代方案是\textbf{SSNDQ 测试}（Small but Significant Non-transitory Decrease in Quality）：假想垄断者能否对用户实施"小幅但显著且非暂时性的质量降级"而仍不流失用户。SSNDQ 把测试的维度从价格转向质量：\important{消费者在免费服务中支付的"价格"是注意力与数据，质量的持续下降相当于隐性的涨价}。\mn{SSNIP 的零价格困境与 SSNDQ 替代：数字市场界定是论述题经典素材，免费的消费者侧没有价格维度，市场界定转向注意力、数据、质量维度。}但质量的多维度与不可观测性使 SSNDQ 难以标准化执行，实务中更多作为分析框架而非操作标准。

多边市场的界定是数字反垄断的另一道难题。美国最高法院在 2018 年俄亥俄州诉美国运通案中确立规则：交易型双边市场的两侧必须一并界定（信用卡网络对商家与持卡人的服务构成一个市场），判决后争议不断；中国 2021 年平台经济指南的立场更灵活：\important{可以根据平台经济的特点，界定一个或多个相关市场}（第 6.4.1 节）。此外，数字市场的\textbf{跨界性}（平台包抄使市场边界模糊，第4章）与\textbf{动态性}（市场支配地位可能随技术变迁快速更替）进一步动摇了静态市场界定的基础。

\subsection{HHI 与勒纳指数}

相关市场划定之后，执法者需要度量市场势力：将价格定在竞争水平之上的能力。测度工具分为两类：\textbf{结构性指标}从市场结构推断势力（HHI），\textbf{行为性指标}直接从定价行为读取势力（勒纳指数）。前者用于事前筛查与集中审查，后者用于事后查处中的势力认定。

\begin{definition}[赫芬达尔—赫希曼指数（Herfindahl–Hirschman Index）]
赫芬达尔—赫希曼指数（HHI）是相关市场内所有企业市场份额平方和的 10000 倍，同时反映企业数量与份额分布：完全垄断时 $HHI = 10000$，$n$ 个等份额企业时 $HHI = 10000/n$。
\end{definition}

其计算公式为
\[
HHI = \sum_{i=1}^{n} s_i^2 \times 10000,
\]
其中 $s_i$ 为各企业市场份额：企业越多、份额越均等，HHI 越低。并购审查中，HHI 的变化 $\Delta HHI$ 是筛选交易竞争关注度的门槛指标，但各版指南的门槛并不相同：美国 1982 年版《横向并购指南》以 $HHI > 1800$ 且 $\Delta HHI > 100$ 作为推定增强市场势力的门槛；现行 2010 年版上调为 $HHI > 2500$ 且 $\Delta HHI > 200$。\mn{各版 HHI 门槛易混：美国 1982 版 1800/100、2010 版 2500/200；中国 2024 年《横向经营者集中审查指引》1000/1800 档位。}中国 2024 年发布的《横向经营者集中审查指引》采用 1000/1800 档位，集中后 HHI 高于 1800 且 $\Delta HHI$ 高于 200 的，通常推定具有排除、限制竞争效果。

\begin{definition}[勒纳指数（Lerner index）]
勒纳指数是价格中超过边际成本的加价比例 $L = (P - MC)/P$，取值介于 0 与 1 之间：$L = 0$ 对应完全竞争定价（$P = MC$），$L$ 越接近 1，市场势力越强。
\end{definition}

勒纳指数的取值由需求价格弹性决定，其推导只需垄断定价的基本一阶条件。

\begin{derivation}{勒纳指数的推导}{}
\textbf{勒纳指数的推导。}设企业面对需求函数 $q(P)$，反需求函数为 $P(q)$，成本函数为 $C(q)$，利润为
\[
\pi(q) = P(q)\, q - C(q).
\]
利润最大化的一阶条件（边际收益等于边际成本）：
\[
\frac{d\pi}{dq} = P + q\, \frac{dP}{dq} - C'(q) = 0.
\]
将边际收益改写为弹性的形式。需求价格弹性 $\varepsilon = \frac{dq}{dP} \frac{P}{q} < 0$，故 $\frac{q}{P} \frac{dP}{dq} = \frac{1}{\varepsilon}$，于是
\[
MR = P + q\, \frac{dP}{dq} = P\left(1 + \frac{q}{P} \frac{dP}{dq}\right) = P\left(1 + \frac{1}{\varepsilon}\right) = P\left(1 - \frac{1}{|\varepsilon|}\right).
\]
代入一阶条件并整理：
\[
P\left(1 - \frac{1}{|\varepsilon|}\right) = MC \quad\Rightarrow\quad L \equiv \frac{P - MC}{P} = \frac{1}{|\varepsilon|}.
\]

勒纳指数与 HHI 之间存在直接的数量联系。同质产品古诺竞争中，企业 $i$ 的一阶条件为 $P + q_i P'(Q) = c_i$，两边同除以 $P$ 并以 $s_i = q_i / Q$ 代入：
\[
\frac{P - c_i}{P} = - \frac{q_i}{P} P'(Q) = - s_i \frac{P'(Q) Q}{P} = \frac{s_i}{|\varepsilon|}.
\]
以市场份额加权的行业平均勒纳指数为
\[
\bar L = \sum_{i=1}^{n} s_i L_i = \frac{\sum_{i=1}^{n} s_i^2}{|\varepsilon|} = \frac{HHI}{10000\, |\varepsilon|}.
\]
两个指标由此贯通：\important{需求弹性给定时，集中度越高，行业平均加价比例越高}——这正是 HHI 得以充当并购筛查工具的微观基础。
\end{derivation}

勒纳指数度量的是\important{定价权，而非利润本身}：$L$ 高说明定价不受竞争约束，但它并不直接等于超额利润。数字市场的适用需注意三点。其一，零边际成本业务（$MC \approx 0$）使 $L \to 1$，形式上"完全垄断"，但若高加价服务于固定成本回收（研发、服务器投入），$L$ 高并不当然意味着超额利润。其二，多边市场的价格是价格结构而非单一价格（第4章），补贴侧的价格低于边际成本时 $L$ 甚至为负，市场势力的读数必须两侧联合判断。其三，动态竞争中短期的高 $L$ 可能是创新租金，测度指标必须与进入壁垒、潜在竞争结合解释。这也解释了实务中为何 HHI 更常用：数字服务的边际成本难以观测，勒纳指数难以直接计算。

\section{平台反垄断的制度与实践}

前三节完成了"垄断形成、垄断行为、测度工具"的理论链条，本节落到制度层面：各国如何组织这些工具、如何回应数字市场的特殊性。先看中国的制度框架与执法范式，再看欧盟的事前监管实验，最后以标志性案例检验分析范式。

\subsection{中国的制度框架}

中国平台反垄断的制度框架成形于 2020—2022 年的监管转折。此前平台经济长期处于"包容审慎"的监管基调之下。转折发生在 2020 年前后：社区团购的补贴混战与"二选一"的蔓延，使平台竞争失序、挤压中小经营者的现象进入公众视野；大数据杀熟与算法歧视的曝光，暴露出数据与算法滥用超出传统监管框架的调整能力；资本以平台为通道无序扩张进入民生领域，则使问题超出市场竞争范畴、触及经济社会秩序。同期，欧盟《数字市场法》进入立法程序、美国执法机构转向积极执法，数字监管成为全球共识。2020 年 12 月中央经济工作会议首提"强化反垄断和防止资本无序扩张"，\mn{中国平台反垄断时间锚点：2021 年 2 月指南、2021 年 4 月阿里案、2022 年 6 月《反垄断法》修订（8 月施行）。}2021 年 2 月《国务院反垄断委员会关于平台经济领域的反垄断指南》发布，平台经济的制度基调从"包容审慎"转向\important{规范发展}。

这一转向的制度支柱有三。其一，\textbf{《反垄断法》}（2022 年修订）：新增第 9 条"经营者不得利用数据和算法、技术、资本优势以及平台规则等从事本法禁止的垄断行为"，将数据、算法纳入垄断行为的认定框架；第 19 条对垄断协议的组织者与实质帮助者设立责任，规制轴辐协议（hub-and-spoke arrangements）；第 26 条允许对未达申报标准但具有或可能具有排除、限制竞争效果的经营者集中要求申报（第 6.2 节数字并购的兜底条款）。其二，\textbf{《国务院反垄断委员会关于平台经济领域的反垄断指南》}（2021）：明确平台经济领域相关市场界定方法（可为多边市场分别界定）、市场支配地位认定因素（网络效应、锁定效应、数据掌握、算法能力）、以及"二选一"、大数据杀熟、"最惠国待遇"条款（MFN 条款）等行为的分析框架。其三，\textbf{执法实践}：阿里巴巴"二选一"案（182.28 亿元）、美团"二选一"案（34.42 亿元）确立了对平台排他性行为的处罚范式。

上述制度在个案中的展开形成一套稳定的执法范式，案例题的分析按五步展开，每一步都有学理依据。第一步\textbf{相关市场界定}：划定竞争约束的范围（第 6.3 节），界定的宽窄直接决定结论，界定过窄会夸大势力、过宽会稀释约束。第二步\textbf{市场支配地位认定}：以市场份额为基础，结合市场进入壁垒、经营者对平台的依赖程度、网络效应与锁定效应、数据与算法能力综合认定；学理上，支配地位是"不受竞争约束地定价"的能力，其认定是滥用行为规制的前提。《反垄断法》第 24 条设有\textbf{支配地位推定}：一个经营者的市场份额达到二分之一、两个经营者合计达到三分之二、三个经营者合计达到四分之三的，可以推定具有市场支配地位，其中份额不足十分之一的经营者不在此列。\mn{支配地位推定数字（2022 年修法第 24 条）：一家 1/2、两家合计 2/3、三家合计 3/4；不足 1/10 的除外；推定可被证据反驳，是选择题与判断题的常考数字。}推定条款的作用是降低执法机关对支配地位的举证负担，经营者可以提交证据反驳推定，因此数字平台的"市场份额超过 50\%"只构成推定的起点而非认定的终点。现代反垄断法以行为规制（conduct-based）而非结构规制（structuralism）为主：\important{支配地位本身不违法，滥用才违法}。第三步\textbf{滥用行为认定}：识别行为类型（排他性协议、自我优待、差别待遇、最惠国待遇条款等，第 6.2 节）。第四步\textbf{竞争损害论证}：行为须产生排除、限制竞争的效果，而非仅损害个别竞争者，反垄断法保护的是竞争机制而非竞争者。第五步\textbf{效率抗辩}：经营者可以证明行为为效率所必需、消费者分享收益、且不严重限制竞争，这是第 6.2 节 Williamson 权衡在举证规则上的制度化，\important{反垄断法的最终标准是福利权衡而非行为的外观}。

\begin{misconception}
\textbf{误区：平台规模大、市场份额高，就应当被拆分或处罚。}反垄断法规制的是\textbf{滥用}市场支配地位的行为，而非支配地位本身。规模大可能是竞争胜出的结果（网络效应下用户用脚投票），也可能是效率的证据。只有当平台利用支配地位实施排除、限制竞争的行为且无效率抗辩时，才构成违法——这正是五步范式中"支配地位"与"滥用行为"分列两步的原因。同理，守门人制度针对的也是特定行为的义务清单，而非拆分公司。
\end{misconception}

\subsection{欧盟《数字市场法》的事前监管}

欧盟《数字市场法》（DMA，2022 年通过、2023 年适用）代表了与事后执法不同的监管范式，其适用对象正是 6.1 节的守门人。2023 年 9 月欧委会指定首批 6 家守门人（Alphabet、亚马逊、苹果、字节跳动、Meta、微软）及其 22 项核心平台服务，义务自 2024 年 3 月起适用。DMA 对守门人施加\textbf{事前义务}：允许商业用户绕过平台接触最终用户（禁止强制使用平台内支付、允许开发者引导用户至平台外支付渠道）、允许用户卸载预装应用、禁止利用平台数据为自营业务谋利（自我优待限制，第 6.2 节）、并购前强制申报。\mn{DMA/DSA 分工口诀：DMA 管行为（竞争），DSA 管内容（治理）。}与 DMA 并列的《数字服务法》（Digital Services Act，DSA）承担平台内容治理职能：要求平台履行内容审核与透明度义务、公开推荐算法的运作方式，并对超大型平台施加风险评估与独立审计义务。DMA 管市场行为、DSA 管内容责任，二者共同构成欧盟数字治理的"两法"体系。事前监管的理论逻辑已在 6.1 节说明：正反馈使损害不可逆、事后救济失灵。其代价是"假阳性"（false positive）风险：义务可能约束了本应合法的商业行为（如平台服务整合的效率），这正是事前与事后两种范式之间经典权衡的数字化重演。

两种范式的权衡可以用决策理论的框架精确表述。事后执法（反垄断查处）的差错成本是\textbf{假阴性}（false negative）：损害已经发生且往往不可逆（临界规模越过、竞争对手退出、数据壁垒固化），即使事后处罚，市场结构已无法修复；事前监管（DMA 义务清单）的差错成本是\textbf{假阳性}（false positive）：义务约束了本无危害的商业行为，抑制平台的投资与创新激励。\important{范式的选择取决于两类误差的相对成本}：当市场自我修复能力弱（网络外部性强、退出成本高）、损害不可逆程度高时，\important{事前监管的假阳性成本低于事后执法的假阴性成本}，事前范式占优——数字市场恰是这一情形。这一定位也划定了事前监管的边界：义务清单应针对"损害不可逆性最强"的行为（自我优待、排他限制），而非一切平台行为，DMA 的守门人门槛（规模+稳固性）正是对"假阳性风险敞口"的限定。对中国而言，2022 年修订的《反垄断法》仍以事后执法为主，但平台监管实践中"合规指引+约谈"的实践已带有事前特征；是否引入成文的事前监管制度，取决于对中国数字市场竞争动态的判读，而判断的标尺正是上述两类误差的比较。

\subsection{标志性案例的分析范式}

制度与工具最终要落到个案。三个标志性案例分别对应三类行为的执法展开：阿里巴巴案对应"二选一"（6.2.3），谷歌购物案对应自我优待（6.2.1），Epic 诉苹果案对应竞争瓶颈收费与市场界定的决定性（第 5 章、第 6.3 节）。案例的价值不在事实本身，而在其展示的分析范式。

\begin{case}[阿里巴巴"二选一"案]
市场监管总局 2021 年 4 月认定：阿里巴巴在中国境内网络零售平台服务市场具有支配地位（市场份额、网络效应、锁定效应、数据与算法能力、资本优势等多因素综合认定）；其通过协议与流量惩罚强制商家"二选一"，剥夺了商家的多归属选择，限制了竞争对手获取商家资源的能力，损害了商家与消费者的利益。处罚：责令停止违法行为，处 2019 年度境内销售额 4\% 的罚款，合计 \important{182.28 亿元}。
\end{case}

\begin{case}[谷歌购物案（欧盟，2017）]
欧盟委员会认定：谷歌在通用搜索服务市场具有支配地位，其在搜索结果中系统性地将自营比价购物服务置于独立比价网站之前（自我优待），导致独立比价网站的流量断崖式下跌，构成滥用市场支配地位。处罚：24.2 亿欧元。该案是全球首例将"自我优待"认定为滥用行为的标志性案件，也是 6.2 节自我优待推导的执法实践。
\end{case}

\begin{case}[Epic 诉苹果案（美国，2021）]
Epic 在游戏《堡垒之夜》中内置自有支付渠道，绕过苹果对应用内数字交易 30\% 的抽成，触发苹果下架与诉讼。2021 年加州北区法院判决：相关市场被界定为"数字移动游戏交易"，苹果在该市场的份额超过 55\% 但未达到垄断势力的认定门槛，30\% 抽成本身未被认定违法；苹果的"反引导"条款（anti-steering，禁止开发者告知用户平台外支付渠道）违反加州《不公平竞争法》，法院发布禁令予以禁止。判决经第九巡回上诉法院维持（2023），美国最高法院拒绝复审（2024）。
\end{case}

Epic 案的范式价值有三层。其一，\textbf{市场界定的决定性}：法院将市场界定为"数字移动游戏交易"而非"iOS 应用分发"，苹果的份额与进入壁垒因而不足以认定垄断势力，抽成水平未进入实体审查，界定的宽窄直接决定了结论（第 6.3 节）。其二，\textbf{竞争瓶颈的收费性质}：应用商店是开发者接触 iOS 用户的唯一渠道（单归属瓶颈，第 4 章），30\% 抽成的可持续性来自锁定而非竞争约束，其上限由被锁定收入的比例决定（第 5 章 $t_{\max} = \ell$）；法院虽未认定垄断，仍以"反引导"条款的救济触及了锁定结构的维护机制。其三，\textbf{事后救济的局限与事前监管的补充}：判决后苹果允许开发者使用外部支付链接，但仍对外部交易收取 27\% 的抽成，抽成结构实质未变——事后个案救济止于边缘行为，这正是 DMA 将"禁止反引导、允许外部支付"直接列为守门人义务（第 6.4 节）的缘由。

\begin{chapterbox}
\textbf{本章考点清单}

\begin{enumerate}
\item 数字市场垄断的三重机制：网络效应 + 规模经济 + 数据反馈回路 \important{【专硕·核心】}
\item 守门人概念与事前监管范式 \important{【专硕·核心】}
\item 自我优待的竞争损害机制、执法先例与规制边界 \important{【专硕·核心】}
\item 杀手型并购、数据驱动型并购、Facebook 收购案与 Williamson 权衡 $\Delta W$ 的推导 \important{【专硕·核心】}
\item "二选一"的竞争瓶颈分析（第4章模型的执法应用）\important{【专硕·核心】}
\item 算法合谋：惩罚时滞条件、四类形态与反垄断责任基础的失灵 \important{【专硕·核心】}
\item SSNIP 测试与临界损失公式 $L_c = x/(x + m/p)$ 的推导 \important{【专硕·核心】}
\item 玻璃纸谬误与多边市场界定（Amex 案）\important{【专硕·核心】}
\item SSNDQ 与零价格市场的界定难题 \important{【专硕·核心】}
\item HHI 与勒纳指数的计算及数字市场适用 \important{【专硕·核心】【学硕·热点】}
\item 中国制度框架（2021 指南、2022 修订、第 24 条支配地位推定）与阿里/美团/谷歌/Epic 案的分析范式 \important{【专硕·核心】}
\item 欧盟 DMA 与 DSA 的分工 \important{【专硕·扩展】}
\item App Store 抽成的竞争瓶颈分析与 Epic 案的市场界定范式 \important{【专硕·扩展】}
\end{enumerate}
\end{chapterbox}

反垄断划定了竞争的行为边界，但数字经济的制度基础不止于此。平台的市场势力、算法的定价能力、监管的执法工具，最终都汇聚于同一个生产要素——数据。数据如何确权、定价并进入资产负债表，是数字经济制度建设的另一半工程。
